\begin{proposition}[DP25: the induced CM type blocks a larger-field substitution]\label{DP25}
Let a CM field $E$ contain an imaginary quadratic field $K$, and let
\[
 \Phi=\{\tau:E\hookrightarrow\mathbf C:\tau|_K=\kappa\}
\]
be the full CM type induced from one embedding $\kappa$ of $K$. Assume good reduction and full integral CM at $p$, with $p$ unramified in $E$. If $p$ is inert in $K$, every rank-one local CM factor has slope $1/2$; a one-dimensional factor has height two. If $p$ splits in $K$, each factor has slope zero or one; a one-dimensional connected factor has height one. Enlarging $E$ while retaining this type cannot produce a one-dimensional factor of height greater than two.

The basic CM blocks of the inspected theta construction have precisely these induced types. The Jacobian of DP22 cannot replace such a block; balancing it with its conjugate also fails the source's stated surjective comparison.
\end{proposition}
\begin{proof}
Restriction of embeddings from $E$ to $K$ commutes with unramified Frobenius. At an inert prime, Frobenius exchanges the two embeddings of $K$, so membership in $\Phi$ alternates around each orbit. Every orbit has even length $d$, with exactly $d/2$ selected embeddings. At a split prime that membership is constant, so the count is zero or $d$. The full-CM height/type formula \cite[\S5, (10*) and Theorem 3]{Tate352} gives height $d$, dimension equal to the count, and slope their ratio. The assertions follow, without requiring an orbit to be stable under complex conjugation.

The source premise is explicit: \path{Formalism/SharpRestriction/RamificationFrobenius/TotientGaloisClasses/ConjEMemHostFIff.lean}, \path{HostRaw.K} and \path{HostRaw.cmType}, define $K=\Q(\sqrt{-d})$ and select exactly one of its embeddings. In \path{Development/AbelianVarietyTheory/Definitions/CMLattices/IsDiffOf.lean}, \path{HasType}, \path{IsTen}, and \path{CMObj.typeP/typeM} impose the two conjugate induced types and their Serre-tensor construction from a quadratic CM elliptic curve. This is a definition of the blocks, not an inference from definiteness of the unitary groups.

For DP22, write $S=\zeta_7+\zeta_7^2+\zeta_7^4$ and $T=\zeta_7^3+\zeta_7^5+\zeta_7^6$. Multiplication gives $S+T=-1$, $ST=2$, and $(S-T)^2=-7$. Thus the unique quadratic subfield has embedding classes $\{1,2,4\}$ and $\{3,5,6\}$. The type $\{1,2,3\}$ is neither class. Relabeling by a Galois automorphism preserves or exchanges its counts $(2,1)$; an $E$-linear isogeny preserves the type. At $11$, $-7\equiv4$ is a square, so that prime splits in $K$. It is outside the source's \path{InertK} gain class, by \path{LocalRingCardinality/DataHost.lean}, \path{primeClass}, and \path{SharpRestriction/GlobalTheoremInputs.lean}, \path{InertK}.

For $J^a\times\overline J^{,b}$ the multiplicities on the two quadratic classes are $(a,a,b)$ and $(a,b,b)$, hence become constant only if $a=b$. For $a=b>0$ the target has type $(a,a)$, but \path{CMObj.onto} also requires a surjection from the corresponding copies of the two induced-type blocks. Equal dimensions would make this an isogeny. Those source blocks are ordinary at $11$, whereas this target has only slopes $1/3$ and $2/3$. Isogeny invariance of Newton slopes rules out that comparison. These statements concern the specified CM action and the existing source data, not every possible mixed-type construction.
\end{proof}
